\scanpage{4}
\diagram{\begin{tikzcd}[ampersand replacement=\&,row sep=1.2em]
G\times X \arrow[r,dashed] \& X\times X\\[-3pt]
\rotatebox{-90}{$\supset$} \& {}\\[-3pt]
U \& {}\\[-3pt]
\rotatebox{-90}{$\supset$} \& {}\\[-3pt]
F \arrow[r] \& X \arrow[uuuu]
\end{tikzcd}}

Opération \uuline{régulière} si $F\mathbin{-} X$ fid. \unsure{plats|plate|plat}

\uline{Prop} Si l'opération est régulière, $F_X$ est un sous-schéma en groupes de $G_X$.

\uline{Cor} $N$ est le plus \unsure{gd|grand} sous-schéma en groupes constant de $F_X$.

Pseudo-opération \inserted{(pseudo)}\uline{transitive} si $G\times X\mathbin{-}X\times X$ \uline{dominant}

(dans le cas d'une vraie opération, cel signifie qu'il y a une \inserted{(unique)} orbite dense)

\uline{Proposition} Si $G$ est pseudo-transitive sur $X$, $X$ est pseudo-isomorphe à un espace homogène.

\uline{Sous-tores de} $\underset{\sim}{\mathrm{Cr}}_{n,k}$ \uuline{Prop} Soit $T$ tore (\unsure{supp|qq} déployé) qui pseudo-opère de façon \deleted{\unsure{régulière}} \deleted{\unsure{effective}} \inserted{fidèle} régulière, \inserted{($X$ irréductible)} Alors $F=e\times X$

\deleted{\unsure{($X$ un|($X$ réd.}}

\uline{Cor} 1 $\dim T\leq\dim X$, égalité $\Longleftrightarrow$ pseudotransitivité

\uline{Cor.} 2 Supposons $\dim T=\dim X$, \inserted{$T$ déployé} Alors $X$ \deleted{$\unsure{\simeq|\approx}$}\inserted{torseur sous}\inserted{$G$-pseudoisomorphisme} $T$

\uline{Cor} 3 $X$ irréductible \inserted{dim $n$}, alors conditions équivalentes

a) $X$ rationnel

b) $\underline{\mathrm{Psaut}}(X)\simeq\underline{\mathrm{Cr}}_{n,k}$

c) $\underline{\mathrm{Psaut}}(x)\supset\unsure{\mathbb G_m^n|G_m^n}$

% archive sheet number: 3
